\documentclass[conference]{IEEEtran}
\IEEEoverridecommandlockouts

\usepackage{cite}
\usepackage{amsmath,amssymb,amsfonts}
\usepackage{graphicx}
\usepackage{textcomp}
\usepackage{xcolor}
\usepackage{booktabs}
\usepackage{url}
\usepackage[hidelinks]{hyperref}

\def\BibTeX{{\rm B\kern-.05em{\sc i\kern-.025em b}\kern-.08em T\kern-.1667em\lower.7ex\hbox{E}\kern-.125emX}}

\begin{document}

\title{Satellite-Based InSAR Ground Deformation Monitoring\\
at the Jharia Coalfield Mine Site Using MineGuard}

\author{
  \IEEEauthorblockN{[Author Names]}
  \IEEEauthorblockA{\textit{[Department / Institution]} \\
  [City, Country] \\
  [email@institution.edu]}
}

\maketitle

% -------------------------------------------------------
\begin{abstract}
This paper presents millimetre-precision ground deformation monitoring of
the Jharia coalfield mine site (23.732\textdegree{}N, 86.439\textdegree{}E),
Jharkhand, India, using the MineGuard InSAR processing pipeline and
Sentinel-1 C-band SAR data acquired from 2017 to 2022.
Ten interferometric pairs across both ascending (ASC) and descending (DSC)
orbital passes were processed with ESA SNAP GPT to derive line-of-sight (LOS)
displacement time series. A 2D ASC+DSC vector decomposition resolved true
vertical and east-west (E-W) displacement components at five user-selected
hotspot locations.
Cumulative vertical subsidence reached \(-380\)\,mm at Hotspot P4 (CRITICAL
risk), while Hotspot P2 exhibited a pronounced eastward lateral shift of
\(+182\)\,mm.
Three of four monitored points were classified as HIGH or CRITICAL risk.
Episodic displacement events detected in 2018 and 2020 suggest discrete
geomechanical failure rather than uniform creep.
An AI-generated geotechnical report, produced via the Claude API, corroborated
the findings and issued site-specific remediation recommendations.
\end{abstract}

\begin{IEEEkeywords}
InSAR, Sentinel-1, ground deformation, mine subsidence, Jharia coalfield,
displacement monitoring, 2D vector decomposition
\end{IEEEkeywords}

% -------------------------------------------------------
\section{Introduction}
The Jharia coalfield in Jharkhand, India, is one of the country's largest
reserves of prime coking coal and has been actively mined for over a century.
Extensive underground extraction causes surface subsidence, slope instability,
and lateral ground creep---hazards that threaten mine infrastructure,
residential structures, and personnel safety \cite{singh2018}.
Conventional geodetic surveys (total stations, levelling) provide only sparse
spatial coverage and are costly to repeat at high temporal frequency.

Interferometric Synthetic Aperture Radar (InSAR) overcomes these limitations
by using the phase difference between repeat-pass SAR acquisitions to derive
spatially dense displacement maps at millimetre-level precision
\cite{massonnet1998, ferretti2001}.
The European Space Agency's (ESA) Sentinel-1 constellation
(C-band, 5.546\,cm wavelength, 6-day revisit) provides a freely accessible
archive of Single Look Complex (SLC) scenes dating back to 2014, making it
especially suitable for long-term mine monitoring.

This work demonstrates the end-to-end capability of \textit{MineGuard}, a
web-based InSAR monitoring platform designed for non-specialist mine site
engineers. MineGuard integrates satellite data acquisition, SNAP GPT
interferometric processing, 2D displacement decomposition, and AI-generated
risk reporting in a single guided workflow.

% -------------------------------------------------------
\section{Study Area and Data}
The study site is the Jharia open-pit and underground mine complex centred
at 23.732109\textdegree{}N, 86.439451\textdegree{}E.
A monitoring bounding box of approximately $4 \times 6$\,km was defined over
the active extraction zone.

Ten Sentinel-1 SLC scene pairs were selected from the Copernicus Data Space
Ecosystem (CDSE) archive \cite{cdse2023}: five ascending-pass (Track ASC)
and five descending-pass (Track DSC) pairs, each spanning approximately one
year (Table~\ref{tab:pairs}).
Temporal baselines of $\sim$12 months and perpendicular baselines below
150\,m were enforced to maintain interferometric coherence above a threshold
of 0.15.

\begin{table}[h]
\caption{Sentinel-1 Interferometric Pairs}
\label{tab:pairs}
\centering
\small
\begin{tabular}{ccc}
\toprule
\textbf{Pair} & \textbf{Master Date} & \textbf{Slave Date} \\
\midrule
ASC 1 / DSC 1 & 2017 & 2018 \\
ASC 2 / DSC 2 & 2018 & 2019 \\
ASC 3 / DSC 3 & 2019 & 2020 \\
ASC 4 / DSC 4 & 2020 & 2021 \\
ASC 5 / DSC 5 & 2021 & 2022 \\
\bottomrule
\end{tabular}
\end{table}

% -------------------------------------------------------
\section{Methodology}
\subsection{InSAR Processing}
Interferograms were generated for all 10 scene pairs using the ESA SNAP GPT
interferometric processing graph, which performs co-registration, interferogram
formation, topographic phase removal (using SRTM 3-arc-second DEM), Goldstein
filtering, and terrain-corrected geocoding \cite{snap2023}.
Phase unwrapping recovered cumulative displacements beyond the single-pair
ambiguity of $\pm$13.9\,mm ($\lambda/4$).

\subsection{2D Displacement Decomposition}
Because Sentinel-1 LOS measurements combine vertical and horizontal
contributions, a dual-orbit decomposition \cite{wright2004} was applied at
each epoch where both ASC and DSC interferograms were available:

\begin{equation}
  \begin{bmatrix} d_\text{V} \\ d_\text{EW} \end{bmatrix}
  =
  \begin{bmatrix}
    \cos\theta_\text{A} & -\sin\theta_\text{A}\sin\phi_\text{A} \\
    \cos\theta_\text{D} & -\sin\theta_\text{D}\sin\phi_\text{D}
  \end{bmatrix}^{-1}
  \begin{bmatrix} d_\text{LOS}^\text{A} \\ d_\text{LOS}^\text{D} \end{bmatrix}
  \label{eq:decomp}
\end{equation}

\noindent where $\theta$ is the incidence angle (ASC: 40.0\textdegree{},
DSC: 35.2\textdegree{}), $\phi$ is the flight heading, and subscripts A/D
denote ascending/descending orbit passes.

\subsection{Hotspot Monitoring}
Five monitoring hotspots (P1--P5) were placed by the user on an interactive
satellite map, constrained within the defined crop area. Geographic coordinates
were mapped to SAR pixel indices for per-pixel phase extraction. Displacement
time series and 3D spatial trajectories were computed for each point.

% -------------------------------------------------------
\section{Results and Discussion}
\subsection{Displacement Time Series}
Fig.~\ref{fig:ew} shows the cumulative east-west displacement per year for
all hotspots derived from 2D ASC+DSC decomposition.
P3 exhibits the largest transient E-W excursion, peaking at approximately
$+$590\,mm westward in 2019--2020 before reversing to $+$113\,mm by 2022.
P4 shows a monotonically increasing westward shift culminating at
$-$449\,mm---the largest E-W displacement observed at the site.
P1 and P2 display moderate eastward trends of $+$22\,mm and $+$63\,mm,
respectively.

\begin{figure}[htbp]
  \centering
  \includegraphics[width=\columnwidth]{results/plot_01_horizontal_displacement.png}
  \caption{Cumulative horizontal (East-West) displacement per year at five
           hotspot locations, derived from 2D ASC+DSC InSAR decomposition,
           Jharia mine site, 2017--2022.}
  \label{fig:ew}
\end{figure}

\subsection{3D Displacement Trajectory}
Fig.~\ref{fig:3d} presents the 3D trajectory of each hotspot through
E-W displacement ($x$-axis), vertical displacement ($y$-axis), and time
($z$-axis, decimal year).
P4 traces a steep downward-westward path, consistent with pit-wall failure
driven by stope collapse. P5 records the largest vertical component
($-$332\,mm) combined with a moderate eastward shift ($+$113\,mm), suggesting
a translational slide mechanism. The non-linear trajectories of P1 and P3
reflect episodic displacement events superimposed on a background creep signal.

\begin{figure}[htbp]
  \centering
  \includegraphics[width=\columnwidth]{results/plot_03_3d_trajectory.png}
  \caption{3D hotspot displacement trajectory ($x$: E-W displacement;
           $y$: vertical displacement; $z$: time in decimal years).
           Labels show cumulative EW and vertical values at 2022.}
  \label{fig:3d}
\end{figure}

\subsection{Risk Assessment}
Table~\ref{tab:risk} summarises cumulative displacements and risk ratings
assigned by the MineGuard AI module (Claude API).
Three of four primary hotspots are rated HIGH or CRITICAL.

\begin{table}[h]
\caption{Cumulative Displacement and Risk Classification (2017--2022)}
\label{tab:risk}
\centering
\small
\begin{tabular}{ccccc}
\toprule
\textbf{Hotspot} & \textbf{Risk} & \textbf{Vert. (mm)} & \textbf{E-W (mm)} \\
\midrule
P4 & \textbf{CRITICAL} & $-379.5$ & $+7.4$  \\
P3 & HIGH              & $-270.0$ & $+34.0$ \\
P1 & HIGH              & $-240.7$ & $+76.0$ \\
P2 & MEDIUM            & $-124.2$ & $+182.2$\\
\bottomrule
\end{tabular}
\end{table}

Two critical episodic events were identified:
(i) In June 2018, P3 and P1 recorded simultaneous large displacements
($+$287\,mm E-W and $-$260\,mm vertical, respectively), indicating a
widespread failure episode likely associated with a mine-wide blasting event
or sudden stope collapse;
(ii) In June 2020, P1 reversed direction with a $-$188\,mm E-W jump,
suggesting renewed westward creep following the 2018 event.
These episodic signatures contrast with simple linear subsidence and
underscore the need for high-cadence monitoring.

% -------------------------------------------------------
\section{Conclusion}
This study confirms active and spatially heterogeneous ground deformation at
the Jharia coalfield mine site using a 5-year Sentinel-1 InSAR time series
processed through the MineGuard platform. Key findings are:

\begin{itemize}
  \item Hotspot P4 experienced a CRITICAL cumulative vertical subsidence of
        $-$380\,mm, requiring immediate geotechnical inspection.
  \item Hotspot P2 shows pronounced eastward lateral drift ($+$182\,mm),
        indicative of active slope translation.
  \item Episodic displacement events in 2018 and 2020 point to discrete
        failure mechanisms rather than uniform consolidation.
  \item The MineGuard AI summary module correctly stratified risk across all
        hotspots and generated actionable remediation guidance.
\end{itemize}

Future work will incorporate the full 6-day Sentinel-1 revisit cycle,
PS-InSAR persistent scatterer analysis, and integration of ground-truth
data from prism targets to validate the satellite-derived displacement
estimates.

% -------------------------------------------------------
\begin{thebibliography}{00}

\bibitem{singh2018}
R.~P. Singh and D.~C. Panigrahi,
``Surface subsidence due to underground coal mining in India: monitoring
and prediction,''
\textit{Int. J. Mining Sci. Technol.}, vol.~28, no.~3, pp.~457--465, 2018.

\bibitem{massonnet1998}
D.~Massonnet and K.~L. Feigl,
``Radar interferometry and its application to changes in the Earth's surface,''
\textit{Rev. Geophys.}, vol.~36, no.~4, pp.~441--500, 1998.

\bibitem{ferretti2001}
A.~Ferretti, C.~Prati, and F.~Rocca,
``Permanent scatterers in SAR interferometry,''
\textit{IEEE Trans. Geosci. Remote Sens.}, vol.~39, no.~1, pp.~8--20, 2001.

\bibitem{wright2004}
T.~J. Wright, B.~E. Parsons, and Z.~Lu,
``Toward mapping surface deformation in three dimensions using InSAR,''
\textit{Geophys. Res. Lett.}, vol.~31, no.~1, 2004.

\bibitem{snap2023}
European Space Agency,
``SNAP --- SeNtinel Application Platform,'' version 9.0,
[Online]. Available: \url{https://step.esa.int/main/toolboxes/snap/}

\bibitem{cdse2023}
European Union / ESA / Eumetsat,
``Copernicus Data Space Ecosystem,''
[Online]. Available: \url{https://dataspace.copernicus.eu/}

\end{thebibliography}

\end{document}
